% ----- CHAUPAI 31 -----

\newpage

\subsection*{\deva{एकत्रिंशत्तम चौपाई} (Thirty-First Chaupai)}
\addcontentsline{toc}{subsection}{\deva{एकत्रिंशत्तम चौपाई} (Thirty-First Chaupai) — \deva{अष्ट सिद्धि...}}

\begin{minipage}[t]{0.55\textwidth}
\vspace{0pt}

{\large \linespread{1.5}\selectfont
\deva{अष्ट सिद्धि नौ निधि के दाता।}\\
\deva{अस बर दीन जानकी माता॥}
\par}

\vspace{0.5em}

{\large \linespread{1.5}\selectfont
\telu{అష్ట సిద్ధి నౌ నిధి కే దాతా।}\\
\telu{అస బర దీన జానకీ మాతా॥}
\par}

\vspace{0.5em}

{\large \linespread{1.3}\selectfont
\textit{aṣṭa siddhi nau nidhi ke dātā |}\\
\textit{asa bara dīna jānakī mātā ||}
\par}
\end{minipage}%
\hfill
\begin{minipage}[t]{0.42\textwidth}
\vspace{0pt}
\begin{tcolorbox}[anchorbox]
\textbf{The Truest Wealth:} The greatest treasure (\textit{Nidhi}) or power (\textit{Siddhi}) in life is not magical abilities or massive bank accounts. True wealth is finding your unique purpose, achieving self-mastery, and building a community of friends and mentors who value you.
\vspace{0.5em}\newline The greatest power or treasure anyone can receive is discovering their true purpose and a community that values them.\newline\null\hfill\href{https://www.valmiki.iitk.ac.in/sloka?field_kanda_tid=5&language=dv&field_sarga_value=68}{\textit{Sundara Kanda, 68:26-29}}
\end{tcolorbox}
\end{minipage}

\vspace{0.5em}
\subsubsection*{\textbf{Visual Rhythm Map:}}
\begin{table}[h!]
\centering
\renewcommand{\arraystretch}{1.2}
\setlength{\tabcolsep}{0pt}
\begin{tabularx}{\textwidth}{|*{16}{>{\centering\arraybackslash\hsize=1.0\hsize}X|}}
\hline
\multicolumn{3}{|>{\centering\arraybackslash\hsize=3.0\hsize}X|}{\textbf{\guru{अ}\laghu{ष्ट}} \par (3)} &
\multicolumn{3}{>{\centering\arraybackslash\hsize=3.0\hsize}X|}{\textbf{\guru{सि}\laghu{द्धि}} \par (3)} &
\multicolumn{2}{>{\centering\arraybackslash\hsize=2.0\hsize}X|}{\textbf{\guru{नौ}} \par (2)} &
\multicolumn{2}{>{\centering\arraybackslash\hsize=2.0\hsize}X|}{\textbf{\laghu{नि}\laghu{धि}} \par (2)} &
\multicolumn{2}{>{\centering\arraybackslash\hsize=2.0\hsize}X|}{\textbf{\guru{के}} \par (2)} &
\multicolumn{4}{>{\centering\arraybackslash\hsize=4.0\hsize}X|}{\textbf{\guru{दा}\guru{ता}} \par (4)} \\
\hline
\end{tabularx}
\vspace{-1.5em}
\end{table}

\begin{table}[h!]
\centering
\renewcommand{\arraystretch}{1.2}
\setlength{\tabcolsep}{0pt}
\begin{tabularx}{\textwidth}{|*{16}{>{\centering\arraybackslash\hsize=1.0\hsize}X|}}
\hline
\multicolumn{2}{|>{\centering\arraybackslash\hsize=2.0\hsize}X|}{\textbf{\laghu{अ}\laghu{स}} \par (2)} &
\multicolumn{2}{>{\centering\arraybackslash\hsize=2.0\hsize}X|}{\textbf{\laghu{ब}\laghu{र}} \par (2)} &
\multicolumn{3}{>{\centering\arraybackslash\hsize=3.0\hsize}X|}{\textbf{\guru{दी}\laghu{न}} \par (3)} &
\multicolumn{5}{>{\centering\arraybackslash\hsize=5.0\hsize}X|}{\textbf{\guru{जा}\laghu{न}\guru{की}} \par (5)} &
\multicolumn{4}{>{\centering\arraybackslash\hsize=4.0\hsize}X|}{\textbf{\guru{मा}\guru{ता}} \par (4)} \\
\hline
\end{tabularx}
\vspace{-1.5em}
\end{table}

\vspace{1em}
\textbf{ \deva{सरल-संस्कृतम्}:}\\
 \deva{भवान् अष्ट-सिद्धीनां नव-निधीनां च दाता अस्ति।}\\
 \deva{एतादृशः वरः माता-जानक्या भवतः कृते दत्तः॥}\\


You are the bestower of the eight supernatural powers and nine divine treasures\\
Such was the boon given to you by Mother Janaki (Sita).


\vspace{1em}
\newpage
\textbf{\deva{प्रतिपदार्थः} (Meaning of each word):}
\begin{table}[h!]
\centering
\renewcommand{\arraystretch}{1.2}
\begin{tabularx}{\textwidth}{@{} l l X X @{}}
\toprule
\textbf{Awadhi} & \textbf{Sanskrit} & \textbf{English} & \textbf{Grammar \& Notes} \\
\midrule
\deva{अष्ट सिद्धि} \glossaryicon / \telu{అష్ట సిద్ధి} / \textit{aṣṭa siddhi}  &  \deva{अष्ट-सिद्धिः}  &  8 Supernatural Powers  & e.g., shrinking/expanding \\
\deva{नौ निधि} \glossaryicon / \telu{నౌ నిధి} / \textit{nau nidhi}  &  \deva{नव-निधिः}  &  9 Divine Treasures  & e.g., endless wealth \\
\deva{के दाता} / \telu{కే దాతా} / \textit{ke dātā}  &  \deva{प्रदाता}  &  Giver / Bestower of  &   \\
\deva{अस बर} / \telu{అస బర} / \textit{asa bara}  &  \deva{एतादृशं वरम्}  &  Such a boon  &   \\
\deva{दीन} / \telu{దీన} / \textit{dīna}  &  \deva{दत्तवती}  &  Gave  & Past tense verb \\
\deva{जानकी माता} / \telu{జానకీ మాతా} / \textit{jānakī mātā}  &  \deva{माता जानकी}  &  Mother Janaki (Sita)  &   \\
\bottomrule
\end{tabularx}
\end{table}

\newpage
